\chapter{Main memory}\label{ch:main-memory}

\whatyoufind{The two memory models behind \code{bus[1]}: a fixed-latency controller for fast
exploration and MCsim's cycle-accurate DDR4 controllers for fidelity; how each is selected,
clocked and stamped; the schedulers MCsim brings; and the two things that went wrong at this
boundary and are worth knowing.}

\section{Two models, one interface}

\begin{figure}[htbp]\centering
\begin{adjustbox}{max width=\textwidth}% The LLC, the memory bus and the two memory models, with their clocks.
\begin{tikzpicture}[font=\sffamily\scriptsize, >=Latex, node distance=5mm and 8mm,
  b/.style={draw=octrule, fill=octmem, rounded corners=2pt, inner sep=5pt, align=center, minimum height=12mm},
  bus/.style={b, fill=octio, draw=octpurple},
  mem/.style={b, fill=white},
  w/.style={<->, thick, octink}, l/.style={font=\sffamily\scriptsize, text=octquiet, align=center}]
  \node[b] (llc) {\textbf{LLC}\\period 100\\\code{GetData} $\rightarrow$ read\\\code{WriteBack} $\rightarrow$ write};
  \node[bus, right=of llc] (bus) {\textbf{bus[1]}\\\code{Point2Point}\\period 50: 2 ticks per core cycle\\request 2, data 5 ticks};
  \node[mem, right=of bus, yshift=10mm] (fx) {\textbf{MainMemoryController}\\period 100, queue 32\\every read: \code{m\_memory\_latency} = 250\\stamps \code{DRAM.ENTER/EXIT}};
  \node[mem, right=of bus, yshift=-10mm] (mc) {\textbf{MCsimInterface} $\rightarrow$ MCsim\\period = the LLC's: one DRAM tick per core cycle\\DDR4-2400, scheduler from \code{mcsim\_scheduler}\\read: callback with the data; write: fire-and-forget};
  \draw[w] (llc) -- (bus);
  \draw[w] (bus.east) -- ++(3mm,0) |- (fx.west);
  \draw[w] (bus.east) -- ++(3mm,0) |- (mc.west);
  \node[l, above=1mm of fx] {\code{main\_memory\_type(s)=MainMemoryController} (default)};
  \node[l, below=1mm of mc] {\code{main\_memory\_type(s)=MCsim -p mcsim\_scheduler(s)=FRFCFS}};
\end{tikzpicture}
\end{adjustbox}
\caption{The LLC, the memory bus and the two memory models with their clocks.}
\label{fig:membus-mcsim}
\end{figure}

The LLC talks to memory through \code{bus[1]}, a point-to-point interconnect: a \act{GetData}
becomes a read, a \act{WriteBack} a write, and the fill comes back the same way
(\cref{fig:membus-mcsim}). Whatever answers on the far side is a \code{CommunicationInterface}
consumer, so the LLC does not know which model it is talking to; \key{main\_memory\_type}
selects it.

\paragraph{\code{MainMemoryController}} (the default) is a fixed-latency model: every read is
served \key{m\_memory\_latency} cycles after it is accepted (250 as shipped), from a queue of
\key{processing\_queue\_size} entries (32), at the cores' clock period. It stamps
\stamp{DRAM.ENTER} on acceptance and \stamp{DRAM.EXIT} when the data leaves, so the report's
\col{DRAM} column is exactly the service and \col{L2-DRAM Bus} the two bus legs
(\cref{ch:request-end-to-end}). It is the right model for design-space exploration where memory
is not the variable, and it is what the sweeps use as the baseline.

\paragraph{\code{MCsimInterface}} wraps MCsim, the extensible DRAM memory-controller simulator
vendored under \file{src/MCsim/}: a DDR4-2400 device with configurable channels, ranks, bank
groups and banks, per-requestor queues, and one scheduler chosen by \key{mcsim\_scheduler} from
the directories under \file{src/MCsim/system/}: FCFS, FR-FCFS (and its batched and capped
variants), BLISS, PAR-BS, MEDUSA, ORP, RTMem, ROC, RankReOrder, ReOrder, Round, PipeCAS, and the
predictable controllers AMC, DCmc, MAG, MCMC, PMC and RROF. Each is a \code{system.ini} that
names the request-queue structure, the command generator and the command scheduler; a new
controller is on average about 130 lines. MCsim stamps no DRAM events: the \col{DRAM} column
becomes the whole \code{bus[1]} round trip and \col{L2-DRAM Bus} the admit-to-leave hand-off.
Writes are fire-and-forget. Selecting it:

\begin{lstlisting}[style=shell]
-p "main_memory_type(s)=MCsim" -p "mcsim_scheduler(s)=FRFCFS"
\end{lstlisting}

MCsim is cycle-accurate and correspondingly slower --- five to ten times on EEMBC, and
prohibitive on the SPLASH giants, which is why the memory axis of the sweeps is EEMBC-only. Use
\code{FRFCFS} unless the scheduler is the subject: the \code{Round} and \code{ORP} schedulers
wedge on the shipped workloads.

\section{Clocks at the boundary}

The memory models are \code{ClockedObj}s like everything else, and both tick at the cores'
rate; what differs is what a tick means on the far side.

\begin{itemize}
\item \code{MainMemoryController} has no clock of its own to speak of: a read is served
  \key{m\_memory\_latency} core cycles after it is accepted. The DRAM is a delay, in core cycles.
\item \code{MCsimInterface} is constructed with the \emph{LLC's} period (one tick per core
  cycle) and tells MCsim that each tick is one cycle of a 2.4\,GHz core
  (\code{setCPUClockSpeed} in \file{src/MCsimInterface.cpp}). MCsim keeps its \emph{own} clock:
  its clock-domain crosser advances the DRAM by one command-clock cycle, $t_{CK}$ of the
  configured device (0.83\,ns for DDR4-2400), for every 0.42\,ns of core time, so the memory
  controller and the DRAM run at the device's real rate relative to the core and MCsim's timing
  parameters (\code{tRCD}, \code{tRP}, \code{tCL}, \ldots) stay in DRAM cycles. A request that
  spends $n$ DRAM cycles in MCsim therefore costs about $2n$ core cycles in the report.
\end{itemize}

The 2.4\,GHz figure is the one assumption at this boundary: the trace-driven core has no
frequency of its own, so the ratio between core and DRAM time is set here, in the bridge, and
is the same for every run.

The simulator is deterministic and its reports are byte-identical on Linux and Windows, MCsim
included; the environment check's reference run (\cref{ch:getting-started}) is an MCsim run for
that reason. \env{MCSIM\_PROBE} logs every request, command and completion with MCsim's own
clock when the DRAM model itself has to be examined (\cref{ch:monitoring}).

\begin{trapbox}[CRLF in the MCsim configuration]
MCsim's \code{.ini} parser rejects CRLF line endings with \emph{Malformed Line N (missing
equals)}. A tree checked out on Windows and copied to Linux has them; the environment check
tests for it.
\end{trapbox}

\begin{wherebox}
\begin{itemize}[nosep,leftmargin=1.2em]
\item \file{src/MainMemoryController.cpp}, \file{configuration/MainMemoryController.csv} --- the fixed-latency model
\item \file{header/MCsimInterface.h}, \file{src/MCsimInterface.cpp} --- the bridge to MCsim
\item \file{src/MCsim/} --- MCsim itself
\item \file{src/MCsim/README.md} --- writing a memory controller
\item \file{src/MCsim/system/<Name>/<Name>.ini} --- one per scheduler
\item \file{src/SystemConfigurations/MultiCoreSystem.cpp} --- where \key{main\_memory\_type} is read
\end{itemize}
\end{wherebox}
